Medical foundation models are reused as frozen feature extractors, zero-shot classifiers and report generators, so it matters to distinguish what their internal representations encode from what the models actually use, independently of any downstream head. This survey reviews work that measures or intervenes on those representations in medical vision encoders, image–text models and multimodal language models, organising the methods by family and by the point along the model at which they read the representation. The evidence across this distinction is uneven. Across \Ncore{} studies, clinical findings, anatomy, demographic attributes and acquisition conditions are decodable from frozen encoders in four to seven modalities, yet controls on probe validity are rare, and no study shows that a model's output actually depends on a named concept. The property recovered most consistently is the site and scanner an image comes from, and whether removing it improves generalisation remains untested. A separate strand compares medical with general pre-training, but none of these comparisons hold data, objective, architecture and scale fixed, so the effect of medical pre-training itself remains unidentified. The survey closes with the experiments that would settle these questions.


